\section*{第 \ref{ch:g8:midpoints} 章——中点与平行线}

\begin{solution}{exo:g8:midpoints:1}
$M$ 和 $N$ 分别是三角形 $RST$ 两边的中点：由
中点定理（\cref{thm:g8:midpoints:direct}），
得 $(MN) \parallel (ST)$，且 $MN = \frac{ST}{2} = 3.5$ cm。
\end{solution}

\begin{solution}{exo:g8:midpoints:2}
$I$ 是中点，且 $(IJ) \parallel (BC)$，由\emph{逆定理}
（\cref{thm:g8:midpoints:converse}），$J$ 是 $[AC]$ 的中点，
故 $AJ = \frac{11}{2} = 5.5$ cm。
\end{solution}

\begin{solution}{exo:g8:midpoints:3}
测量得 $IJ = 3.5$ cm $= \frac{BC}{2}$。三角形 $AIJ$ 的
\omterm{def:g6:measure:perimeter}{周长}：$AI = 4$，$AJ = 3$，$IJ = 3.5$：共 $10.5$ cm，恰为
$ABC$ \omterm{def:g6:measure:perimeter}{周长}的\omterm{def:g3:division:half}{一半}（$8 + 6 + 7 = 21$ cm）。
\end{solution}

\begin{solution}{exo:g8:midpoints:4}
$IJK$ 的每条边都是 $ABC$ 对应边的\omterm{def:g3:division:half}{一半}：分别为 $5$、$6$ 和 $8$ cm。
\omterm{def:g6:measure:perimeter}{周长}为 $19$ cm，恰为 $10 + 12 + 16 = 38$ cm 的\omterm{def:g3:division:half}{一半}。
\end{solution}

\begin{solution}{exo:g8:midpoints:5}
中点三角形是
\cref{prop:g8:midpoints:medial}中的四个全等三角形之一，\omterm{def:g6:measure:area}{面积}为 $\frac{36}{4} = 9$ cm$^2$——
这四个小三角形的\omterm{def:g6:measure:area}{面积}都是 $9$ cm$^2$。
\end{solution}

\begin{solution}{exo:g8:midpoints:6}
由 $ABC$ 中的中点定理得：$(IJ) \parallel (BC)$，即
$(IJ) \parallel (BK)$，且 $IJ = \frac{BC}{2} = BK$（因 $K$ 为
$[BC]$ 的中点）。$IJKB$ 有一组对边\omterm{def:g6:lines:perp}{平行}且
相等，所以它是\omterm{def:g7:central:parallelogram}{平行四边形}
（\cref{thm:g7:central:recognize}，判定条件 3）。
\end{solution}

\begin{solution}{exo:g8:midpoints:7}
接缝连接两个中点：$\frac{6.4}{2} = 3.2$ m。在
距\omterm{def:g5:solids:def}{顶点}\omterm{def:g3:division:half}{四分之一}处，图示表明接缝长为
$\frac{6.4}{4} = 1.6$ m，与底边\omterm{def:g6:lines:perp}{平行}——这一结论
在\cref{ch:g9:thales}中由泰勒斯予以证实。
\end{solution}

\begin{solution}{exo:g8:midpoints:8}
\emph{1.} $M$ 是 $[BC]$ 的中点，且 $(MN) \parallel (AB)$：
由中点定理的逆定理（在三角形 $ABC$ 中，
边 $[AC]$），$N$ 是 $[AC]$ 的中点。

\emph{2.} $(MN) \parallel (AB)$，且 $(AB) \perp (AC)$（\omterm{def:g3:shapes:rightangle}{直角}
位于 $A$ 处）：\omterm{def:g6:lines:perp}{平行}于 $(AB)$ 的直线也\omterm{def:g6:lines:perp}{垂直}于 $(AC)$
（\cref{prop:g6:lines:facts}）。故 $(MN) \perp (AC)$——\omterm{def:g6:lines:objects}{线段}
$[MN]$ 是\omterm{def:g6:symmetry:bisector}{垂直平分线}的一部分，再次说明
$MA = MC$（\cref{thm:g8:cosine:circle}）。
\end{solution}

\begin{solution}{exo:g8:midpoints:9}
\emph{1.} 在三角形 $ABC$ 中，$I$、$J$ 分别是 $[AB]$ 和
$[BC]$ 的中点：$(IJ) \parallel (AC)$，且 $IJ = \frac{AC}{2}$。在三角形
$ACD$ 中，$L$、$K$ 分别是 $[DA]$、$[CD]$ 的中点：
$(LK) \parallel (AC)$，且 $LK = \frac{AC}{2}$。

\emph{2.} 故 $[IJ]$ 与 $[LK]$ \omterm{def:g6:lines:perp}{平行}（均与
对角线 $(AC)$ \omterm{def:g6:lines:perp}{平行}），且相等（都是 $AC$ 的\omterm{def:g3:division:half}{一半}）：$IJKL$ 是
\omterm{def:g7:central:parallelogram}{平行四边形}（\cref{thm:g7:central:recognize}，判定条件 3）——
对于\emph{任意}\omterm{def:g4:geometry:polygon}{四边形} $ABCD$，
正如\cref{exo:g7:central:11}中所猜测的那样。
\end{solution}

\begin{solution}{pb:g8:midpoints:1}
\textbf{1.} 无论画出怎样的三角形，三条中线看起来
都交于同一点，该点位于三角形内部，
并且明显地更靠近每条边的中点，而不是
与之相对的\omterm{def:g5:solids:def}{顶点}。本题的其余部分将证明这一点。

\textbf{2.} 取 $[AK]$ 为三角形 $ABC$ 的中线，则
三角形 $ABK$ 与 $AKC$ 的底边 $BK$ 与 $KC$
等长（因 $K$ 是 $[BC]$ 的中点），且都位于同一直线
$(BC)$ 上，并且高也相同，即点 $A$ 到 $(BC)$ 的距离。
由\omterm{def:g6:measure:area}{面积}公式，每个三角形的\omterm{def:g6:measure:area}{面积}都等于
底 $\times$ 高 $\div\ 2$，两者的底与高分别相等，
因此\omterm{def:g6:measure:area}{面积}相等。对于\omterm{def:g6:measure:area}{面积}为
$36$ cm$^2$ 的三角形，两部分的\omterm{def:g6:measure:area}{面积}各为 $18$ cm$^2$。

\textbf{3.} $I$、$J$ 分别是三角形的边 $[AB]$ 和
$[AC]$ 的中点（三角形 $ABC$），由
\cref{thm:g8:midpoints:direct}可得：$(IJ) \parallel (BC)$，且
$IJ = \frac{BC}{2}$。

\textbf{4.} $M$、$N$ 分别是三角形的边 $[GB]$ 和
$[GC]$ 的中点（三角形 $GBC$），由同一定理可得：
$(MN) \parallel (BC)$，且 $MN = \frac{BC}{2}$。

\textbf{5.} $(IJ)$ 与 $(MN)$ 都与 $(BC)$ \omterm{def:g6:lines:perp}{平行}，故
彼此\omterm{def:g6:lines:perp}{平行}；且 $IJ = \frac{BC}{2} = MN$。
\omterm{def:g4:geometry:polygon}{四边形} $IJNM$ 因此有一组对边 $[IJ]$ 与
$[NM]$ \omterm{def:g6:lines:perp}{平行}且等长：由
\cref{thm:g7:central:recognize}的判定条件 3 可知，$IJNM$ 是\omterm{def:g7:central:parallelogram}{平行四边形}。

\textbf{6.} $G$ 位于中线 $[CI]$ 上，且 $N$ 是
$[GC]$ 的中点，它在同一中线所在的直线上：
故 $I$、$N$（以及 $G$、$C$）都位于 $(CI)$ 上。同理，$J$、$M$
及 $G$ 均位于中线 $(BJ)$ 上。在\omterm{def:g7:central:parallelogram}{平行四边形}
$IJNM$ 中，对角线为 $[IN]$ 与 $[JM]$，它们分别位于
两条中线上，两条中线相交于 $G$ ——因此对角线
相交于 $G$。由\omterm{def:g7:central:parallelogram}{平行四边形}的\emph{定义}
（\cref{thm:g7:central:recognize}的判定条件 1：对角线
互相平分），$G$ 是 $[IN]$ 的中点，\emph{并且}
也是 $[JM]$ 的中点：$GI = GN$，且 $GJ = GM$。

\textbf{7.} $M$ 是 $[GB]$ 的中点，故 $BM = MG$；且
由第 6 题可得 $MG = GJ$。于是 $BM = MG = GJ$，中线
$[BJ]$ 被三等分，其中 $BG$ 占两份：
\[
BG = 2 \times GJ .
\]
同理，$N$ 是 $[GC]$ 的中点，且 $GN = GI$，故
$CN = NG = GI$，且 $CG = 2 \times GI$：在中线 $[CI]$ 上同样
交点距\omterm{def:g5:solids:def}{顶点}恰为三分之二处。

\textbf{8.} 上述论证并未用到数对
$[BJ]$、$[CI]$ 的任何特殊之处：若取数对 $[BJ]$、$[AK]$ 重复该论证，
则可知它们的交点也位于 $[BJ]$ 上，且
距 $B$ 恰为三分之二处。然而\omterm{def:g6:lines:objects}{线段} $[BJ]$ 上只有\emph{一个}
点，距离为 $\frac23 BJ$ 且自 $B$ 量起，故
$[BJ]$ 与 $[AK]$ 的交点恰为同一点 $G$。三条中线
因此都经过 $G$，即它们\emph{共点}。

\textbf{9.} $BG = \frac23 \times 9 = 6$ cm，且
$GJ = \frac13 \times 9 = 3$ cm。

\textbf{10.} 由第 8 题，同样的三分之二计算
对中线 $[AK]$ 也成立（以数对 $[AK]$、$[BJ]$ 重复第 4 至 7 题的论证），
即 $AG = 2 \times GK$。

\textbf{11.} 在三角形 $GBC$ 中，点 $K$ 是
边 $[BC]$ 的中点，故 $[GK]$ 是 $GBC$ 的中线；由
第 2 题，三角形 $GBK$ 与 $GKC$ \omterm{def:g6:measure:area}{面积}相等。

\textbf{12.} 中线 $[AK]$ 把 $ABC$ 分成 $ABK$ 与 $AKC$，两者
\omterm{def:g6:measure:area}{面积}相等（第 2 题）。从两者中各去掉
第 11 题所得的小三角形：
\[
\text{面积}(ABG) = \text{面积}(ABK) - \text{面积}(GBK),
\qquad
\text{面积}(ACG) = \text{面积}(AKC) - \text{面积}(GKC),
\]
而两块被去掉的\omterm{def:g6:measure:area}{面积}相等，故
$\text{面积}(ABG) = \text{面积}(ACG)$。同样的计算
沿中线 $[BJ]$——它把 $ABC$ 分成 $ABJ$ 与
$CBJ$，而 $[GJ]$ 是三角形 $GAC$ 的中线，将其分成
$GAJ$ 与 $GCJ$——由此得
$\text{面积}(ABG) = \text{面积}(CBG)$。于是三角形
$ABG$、$BCG$、$CAG$ \omterm{def:g6:measure:area}{面积}相等，且三者合起来
恰好拼成 $ABC$，故每块\omterm{def:g6:measure:area}{面积}为整体的三分之一。

\textbf{13.} 在三角形 $ABG$ 中，点 $I$ 是
边 $[AB]$ 的中点，故 $[GI]$ 是 $ABG$ 的\emph{从\omterm{def:g5:solids:def}{顶点}
$G$ 引出的中线}：它把 $ABG$ 分成两个\omterm{def:g6:measure:area}{面积}相等的三角形
（第 2 题），即 $AIG$ 与 $IBG$。同理，在 $BCG$ 中
（中线 $[GK]$）与 $CAG$ 中（中线 $[GJ]$）情形相同。
$ABC$ 的每个三分之一都被平分，围绕 $G$ 的六块的\omterm{def:g6:measure:area}{面积}各为
$\frac13 \times \frac12 = \frac16$，即整个三角形\omterm{def:g6:measure:area}{面积}的六分之一。

\textbf{14.} 在中线 $[AK]$ 上，从 $A(0,0)$ 到 $K(3,3)$：$G$
位于三分之二处，故其坐标为
$\left(\frac23 \times 3,\ \frac23 \times 3\right) = (2, 2)$。在
中线 $[BJ]$ 上，从 $B(6, 0)$ 到 $J(0, 3)$：三分之二处的
位置等于加上位移的三分之二（水平方向 $-6$，
竖直方向 $+3$）：
\[
\left(6 + \tfrac23 \times (-6),\ 0 + \tfrac23 \times 3\right)
= (2, 2) .
\]
两者为同一点——正如第 8 题所预言。并且
$\frac{0 + 6 + 0}{3} = 2$，$\frac{0 + 0 + 6}{3} = 2$：重心
的坐标恰为三个\omterm{def:g5:solids:def}{顶点}坐标的
\emph{平均数}。

\textbf{15.} 由第 2 题，沿中线放置的刀口两侧
硬纸板的量相同（即\omterm{def:g6:measure:area}{面积}相等），因此
三角形沿三条中线所在的每一条直线都能平衡。平衡点
必须同时位于三条平衡线上，由第 8 题可知，
三条中线恰有一个公共点 $G$。这就是
铅笔尖应放的位置。（两侧\omterm{def:g6:measure:area}{面积}相等是较为可信的论证版本；
较严谨的版本则需衡量
\emph{硬纸板放在何处}，而不仅是其多少，
则需借助本科分册中的积分学——该方法
证实了 $G$ 的位置。）
\end{solution}
